	\begin{center}
		\section{基础算法}
	\end{center}

	\begin{multicols*}{2}

		\subsection{二维前缀和}

		二维前缀和递推公式。

		\begin{minted}{cpp}
Sum[i][j]=Sum[i-1][j]+Sum[i][j-1]+(G[i][j]=='g')-Sum[i-1][j-1];
		\end{minted}

		\begin{minted}{cpp}
ll lans=Sum[xr][yr]+Sum[xl-1][yl-1]-Sum[xl-1][yr]-Sum[xr][yl-1];
		\end{minted}

		或者调用库函数计算二维前缀和。

		\begin{minted}{cpp}
class NumMatrix {
public:
    vector<vll> Sum;
    ll N,M;
    NumMatrix(vector<vector<int>>& matrix) {
        N=SZ(matrix),M=SZ(matrix[0]);
        Sum.assign(N+5,vll(M+5,0));
        FOR(i,0,N-1){
            partial_sum(all(matrix[i]),Sum[i+1].begin()+1);
        }
        FOR(i,2,N){
            FOR(j,1,M){
                Sum[i][j]=Sum[i][j]+Sum[i-1][j];
            }
        }
    }
    int sumRegion(int row1, int col1, int row2, int col2) {
        row2++;col2++;
        ll ans=Sum[row2][col2]+Sum[row1][col1]-Sum[row1][col2]-Sum[row2][col1];
        return ans;
    }
};
		\end{minted}

		\subsection{二维差分}

		\begin{minted}{cpp}
class Solution {
public:
    vector<vector<int>> rangeAddQueries(int n, vector<vector<int>>& queries) {
        vector<vii> diff(n,vii(n));
        FOR(i,0,SZ(queries)-1){
            ll l1=queries[i][0],r1=queries[i][1],l2=queries[i][2],
            r2=queries[i][3];
            diff[l1][r1]+=1;
            if(r2+1<n) diff[l1][r2+1]-=1;
            if(l2+1<n)diff[l2+1][r1]-=1;
            if(l2+1<n && r2+1<n) diff[l2+1][r2+1]+=1;
        }
        FOR(i,0,n-1){
            partial_sum(all(diff[i]),diff[i].begin());
        }
        FOR(i,1,n-1){
            FOR(j,0,n-1){
                diff[i][j]=diff[i][j]+diff[i-1][j];
            }
        }
        return diff;
    }
};
		\end{minted}

		\subsection{双指针模板}

		\textbf{参考来源}：\href{https://www.acwing.com/}{AcWing}

		\begin{minted}{cpp}
for (int l = 0, r = 0; r < n; r++) {
    add(a[r]);    // 1 扩展：a[r] 进入窗口
    while (l < r && check(l, r)) l++; // 2 收缩
    // 更新答案    // 3 记录
}
		\end{minted}

		\textbf{三段顺序的硬约束}：

		\begin{enumerate}
			\item 步骤 1 \textbf{必须}在步骤 2 之前——\texttt{check(l, r)} 要用到 \texttt{a[r]} 已在窗口的事实。
			\item 步骤 3 \textbf{必须}在步骤 2 之后——只有收缩到位的 $[l, r]$ 才是当前 $r$ 下的最优左端。
		\end{enumerate}

		AcWing 原始模板把扩展和记录都塞进「具体问题的逻辑」一行注释里，容易让人误以为收缩在记录前面就够了；实际上扩展、收缩、记录三步都不能省、相对位置也不能换。

		\textbf{常见问题分类}：

		\begin{enumerate}
			\item 对于一个序列，用两个指针维护一段区间。
			\item 对于两个序列，维护某种次序，比如归并排序中合并两个有序序列的操作。
		\end{enumerate}

		\subsection{整数二分}

		\begin{minted}{cpp}
bool check(int x) {/* ... */} // 检查x是否满足某种性质

// 区间[l, r]被划分成[l, mid]和[mid + 1, r]时使用：
int bsearch_1(int l, int r)
{
    while (l < r)
    {
        int mid = l + r >> 1;
        if (check(mid)) r = mid;    // check()判断mid是否满足性质
        else l = mid + 1;
    }
    return l;
}
// 区间[l, r]被划分成[l, mid - 1]和[mid, r]时使用：
int bsearch_2(int l, int r)
{
    while (l < r)
    {
        int mid = l + r + 1 >> 1;
        if (check(mid)) l = mid;
        else r = mid - 1;
    }
    return l;
}
		\end{minted}

		也可以用「两端哨兵」写法，核心是先固定不变量：一个端点始终可行，另一个端点始终不可行。这样最后返回哪边不会混。

		\begin{minted}{cpp}
// 最大化答案：check(x) 为 T,T,T,F,F
// invariant: check(lo) == true, check(hi) == false
int lo = L - 1, hi = R + 1;
while (hi - lo > 1) {
    int mid = (lo + hi) / 2;
    if (check(mid)) lo = mid;
    else hi = mid;
}
return lo;

// 最小化答案：check(x) 为 F,F,F,T,T
// invariant: check(lo) == false, check(hi) == true
int lo = L - 1, hi = R + 1;
while (hi - lo > 1) {
    int mid = (lo + hi) / 2;
    if (check(mid)) hi = mid;
    else lo = mid;
}
return hi;
		\end{minted}

		\subsubsection{partition\_point 实现二分答案（C++20）}

		可以知道 \texttt{T,T,T,T,F,F,F,F} 中的第一个 \texttt{F} 出现位置，\texttt{iota} 实现的是这个左闭右开区间，如果想反过来，就是像 \texttt{range2} 一样使用这个 \texttt{reverse} 视图即可。

		\begin{minted}{cpp}
auto check_bs=[&](ll i) {
    return query_is_a_gt_b(i,sec_down_start,memo);
};
auto range1=views::iota(sec_down_start+1,end_of_down);
auto it1=ranges::partition_point(range1,check_bs);
auto range2=views::iota(end_of_down,N+1)|views::reverse;
auto it2=ranges::partition_point(range2,check_bs);

it1=ranges::prev(it1);

ll to_be_select=*it1;
if (it2!=range2.begin()) {
    it2=ranges::prev(it2);
    if (query_is_a_lt_b(*it2,*it1,memo)) {
        to_be_select=*it2;
    }
}
		\end{minted}

		\subsection{三分（实数）}

		凸函数。

		\begin{minted}{cpp}
while(R-L>eps){    // 保证为凸性函数在[l,r]内
    DB mid=L+(R-L)/2;
    if(F(mid-eps)>F(mid+eps)){
        R=mid;
    }else{
        L=mid;
    }
}
		\end{minted}

		还有一种实践方式：使用固定迭代次数代替 \texttt{eps} 阈值。

		\begin{minted}{cpp}
// 方案二：使用固定迭代次数（100次）进行三分搜索。
// 这种方法不依赖eps，更稳定，且100次迭代足以保证精度远高于1e-8。
for(int _=0; _<100; ++_){   // 凹函数，求最小值用的
    DB mid1 = l + (r-l)/3.0; // 推荐用 l + (r-l)/3 的形式，数值上更稳定
    DB mid2 = r - (r-l)/3.0;
    if(fx(mid1) > fx(mid2)){
        l=mid1;
    }else{
        r=mid2;
    }
}
printf("%.11f\n",fx(l));
		\end{minted}

		\subsection{三分整数}

		注意\textbf{整数三分遇到平台都会出现问题}，调整 \texttt{r-l>?} 这个判定阈值也只是缓兵之计，只能保证在小于这个平台的数据上得到正确答案。

		因此正确的想法应该是消除平台（尽可能），比如设计这个 $f$ 函数，\textbf{使其返回值是一个浮点数}，但是仍能反映增减关系（如 $f$ 原来是\textbf{向下取整除法}，可改成\textbf{浮点数除法}）。

		\begin{minted}{cpp}
l=0,r=N-1;
while (r-l>2) {
    ll mid1=l+(r-l)/3;
    ll mid2=r-(r-l)/3;
    if (f(mid1)<=f(mid2)) {
        l=mid1;
    }else {
        r=mid2;
    }
}
FOR(i,l,r) ans=max(ans,f(i));
		\end{minted}

		参考提交：\href{https://atcoder.jp/contests/abc421/submissions/68978595}{AtCoder ABC 421}。

		另一种处理思路：将 $f$ 函数反转（比如加一个符号）使凸函数变为凹函数（相当于沿 X 轴翻折下来），再用普通二分定位极值。

		\begin{minted}{cpp}
ll l=1,r=mn;
while (l<r) {
    ll mid=l+r>>1;
    // if(f(mid)<=f(mid+1)){
    // 从mid点（或其前）开始，函数开始单调递增
    if(f(mid+1)-f(mid)>=0){
        r=mid;
    }else{
        l=mid+1;
    }
}
		\end{minted}

		\subsection{倍增法}

		参考题目：\href{https://www.matiji.net/exam/oj/submit-detail/13507760/29/4498/F16DA07A4D99E21DFFEF46BD18FF68AD}{Matiji 提交记录}。

		首先求出 \texttt{nextPos} 数组（注意：该数组定义是\textbf{结束位置}，所以 \texttt{break} 以后需要 \texttt{-1}），然后从 \texttt{N+1} 开始往前递推。

		\begin{minted}{cpp}
vector<vector<ll> > jump(N+5,vector<ll>(mx_k+2,0));
ROF(i,N+1,1){
    jump[i][0]=i<=N?nextPos[i]+1:N+1;
    FOR(j,1,mx_k){
        jump[i][j]=jump[jump[i][j-1]][j-1];
    }
}
		\end{minted}

		调用方式：

		\begin{minted}{cpp}
FOR(i,1,N){
    ll lans=0,idx=i;
    FOR(j,0,mx_k){
        if(bik[j]){
            idx=jump[idx][j];
        }
    }
    lans=idx-i;
    ans=max(ans,lans);
}
		\end{minted}

		\subsection{LIS 最长上升子序列及其还原}

		\textbf{0-based 输入}，构造时传 \texttt{strict} 开关：\texttt{true} 求严格递增 LIS，\texttt{false} 求单调不降 LNDS（最长非降子序列）。算法核心是 patience sort：严格用 \texttt{lower\_bound} 替换首个 $\geq a_i$，非严格用 \texttt{upper\_bound} 替换首个 $> a_i$。还原阶段反向贪心，严格判 \texttt{<}，非严格判 \texttt{<=}。

		参考提交：\href{https://judge.yosupo.jp/submission/327571}{Yosupo Library Checker}（严格）、\href{https://judge.yosupo.jp/submission/370233}{OpenJudges 单调不降}（长度，非严格还原已通过本地大规模对拍 \(O(N^2)\) 暴力 DP 验证 512/512）。

		\begin{minted}{cpp}
// 最长 (严格 / 非严格) 上升子序列 + 还原值与下标
// 0-based 输入；strict=true → 严格递增, strict=false → 单调不降
struct LIS {
    vector<ll> a;       // 输入序列
    int n;
    bool strict;
    vector<ll> tail;    // patience sort: 长度-i 的子序列最小末值
    vector<int> i_len;  // 以下标 i 结尾的 LIS 长度
    vector<ll> seq;     // 还原的递增子序列值
    vector<int> pos;    // 还原的下标 (0-based, 严格递增)

    LIS(const vector<ll>& v, bool strict_ = true)
        : a(v), n((int)v.size()), strict(strict_), i_len(v.size(), 0) {
        for (int i = 0; i < n; ++i) {
            // 严格 → lower_bound（替换首个 >= a[i]）
            // 非严格 → upper_bound（替换首个 > a[i]）
            auto it = strict
                ? lower_bound(tail.begin(), tail.end(), a[i])
                : upper_bound(tail.begin(), tail.end(), a[i]);
            if (it == tail.end()) {
                tail.push_back(a[i]);
                i_len[i] = (int)tail.size();
            } else {
                int idx = (int)(it - tail.begin());
                tail[idx] = a[i];
                i_len[i] = idx + 1;
            }
        }
        // 反向贪心还原一组解
        int need = (int)tail.size();
        ll cur = LLONG_MAX;
        for (int i = n - 1; i >= 0 && need > 0; --i) {
            bool ok = strict ? (a[i] < cur) : (a[i] <= cur);
            if (ok && i_len[i] >= need) {
                seq.push_back(a[i]);
                pos.push_back(i);
                cur = a[i];
                --need;
            }
        }
        reverse(seq.begin(), seq.end());
        reverse(pos.begin(), pos.end());
    }
    int size() const { return (int)tail.size(); }
};

void Solve(){
    int N; cin >> N;
    vector<ll> P(N);
    for (int i = 0; i < N; ++i) cin >> P[i];
    LIS lis(P);                 // 默认严格；非严格写 LIS lis(P, false);
    cout << lis.size() << "\n";
    for (int p : lis.pos) cout << p << " ";
    cout << "\n";
}
		\end{minted}

		\subsubsection{Dilworth 定理（最小链覆盖 / 偏序划分）}

		偏序集 $(S,\preceq)$ 上，两两\emph{可比}的子集叫\textcolor{blue}{\textbf{链}}，两两\emph{不可比}的子集叫\textcolor{blue}{\textbf{反链}}。

		\textbf{Dilworth 定理}：把 $S$ 划分成链的\emph{最少链数} $=$ \emph{最长反链}的元素个数。\textbf{对偶（Mirsky 定理）}：把 $S$ 划分成反链的\emph{最少反链数} $=$ \emph{最长链}的长度。

		\textbf{落到序列上}（与上面的 \texttt{LIS} 互通），下标先后即时间序：

		\begin{itemize}
			\item 划分成\textbf{最少条不升}子序列 $=$ 最长\textbf{严格上升}子序列长度；
			\item 划分成\textbf{最少条不降}子序列 $=$ 最长\textbf{严格下降}子序列长度；
			\item 对偶：最长\textbf{不降}子序列长度 $=$ 划分成\textbf{最少条严格下降}子序列的数目。
		\end{itemize}

		\textbf{用法}：说白了，从``最少链覆盖''翻到``最长子序列''，只是同时拨两个开关——\textcolor{blue}{\textbf{严格}}$\leftrightarrow$\textcolor{blue}{\textbf{非严格}}、\textcolor{blue}{\textbf{上升}}$\leftrightarrow$\textcolor{blue}{\textbf{下降}}，再套上面那个 \texttt{LIS} 即可。

	\end{multicols*}

